\scanpage{14}
% L1
conclurait que $\tau_i$ est dans le centre -- ce qui
% L2
n'est pas le cas pour \inserted{ni pour} $i=0$ ni $1$ \deleted{\unsure{ou}} $2$, si $(\tau_0,\tau_1)$ et
% L3
$(\tau_1,\tau_2)$ ne commutent pas. Dans ce cas, on a
% L4
donc $\dot\tau_0,\dot\tau_1,\dot\tau_2$ sont égaux à l'unique élément $\ne1$
% L5
dans $\unsure{\gamma|\mathcal T}$, donc $\rho_0$ et $\rho_1$ sont dans $H$, donc
% L6
commutent. \unsure{$=$}

% L7
$C\Rightarrow\deleted{\unsure{D|A}}C'$) C implique que le \deleted{image} gpe engendré
% L8
par les $\tau_i$ est isom à un quotient d'un
% L9
des trois gpes déduits de $G_2$ en
% L10
divisant par l'une des relations $\deleted{(\tau_0\tau_1)^2=1}$ $\rho_0^2=1$
% L11
$\deleted{\tau_1\tau_2}\ \underline{\rho}_1^{\deleted{\unsure{0}}2}=1$ ou $\underline{\rho}_0\underline{\rho}_1=\rho_1\rho_0$. On voit \unsure{tt|tout} de
% L12
suite que chacun de ces groupes est
% L13
isomorphe à $\Gamma'=\mathbb Z_{/2}\times\mathbb D_\infty$,

% L13a
$C'\Rightarrow C''$\quad \deleted{\unsure{Il suffit que} \unsure{de}}\inserted{\unsure{qui contient}}
% L14
\deleted{un tel groupe, on a} un ss groupe invariant abélien \inserted{cyclique} $\mathbb D_\infty^+=\Gamma''\subset\Gamma'$, 
% L15
à quotient abélien

% L16
$C'')\Rightarrow A)$ Si $u$ est un gén\unsure{,} de $\Gamma''$, et si $u_{\unsure{Y|\gamma|\varphi}}=\mathrm{id}$, l'hom
% L17
se factorise par le gpe comm. \uline{quotient}, et on
% L18
gagne (....). Sinon, le diviseur des pts fixes de
% L18a
\margin{\uline{réf?}}
% L19
$u_Y$ est intrinsèquemt attaché \inserted{\unsure{donc|sous}} au \deleted{\unsure{$q$}} gpe des $\unsure{\mathrm{Aut}_k(Y)|\mathrm{Aut}_Y(k)|\mathrm{Aut}(X)}$
% L20
engendré par $u_Y$, donc stable par $\Gamma'$, on gagne.

% L21
Donc A) à C'') sont équivalents, d'ailleurs
% L22
$D\Rightarrow C'')$, montrons que B) \& C') implique
% L23
D.... Il suffit de prouver que tt ss-gpe d'
% L24
\deleted{$B$ ou $N(T)$} d'un $H=B$ ou $N(T)$ qui est $\inserted{\simeq}$ un
% L25
quotient de $\mathbb D_\infty\times\mathbb Z_{/2}$, est isom. à $\mathbb D_\infty\times\mathbb Z_{\unsure{2|/2}}$ ou
% L26
ou à un $\mathbb D_n$, ou à un $\mathbb D_n\times\mathbb Z_2$.
% archive sheet number: 13
